The dataset used in this thesis is a subset of a larger dataset obtained from a company in Hamburg.  The dataset is a collection of aerial images captured by manned aerial vehicles. The images were obtained at different times of the day and year and under varying weather conditions, resulting in a diverse set of environmental and lighting conditions, e.g sometimes the water is blue, sometimes gray. 

\begin{figure}[!h]
	 \centerline{\includegraphics[width=1\textwidth]{src/approach/distribution_dataset.png}}
	 {\caption{Distribution of the dataset} \label{fig:dataset_distribution}}
\end{figure}

The dataset comprises more than 2,700 images with a resolution of 500 x 500 pixels. Ornithologist annotated the images to identify the marine wildlife and anthropological artifacts present in each image. The objects of interest include different species of marine mammals, sea turtles, and various bird species. Illustrated in Figure \ref{fig:dataset_distribution}, the dataset consists of 118 classes and is divided into four primary types: anthropological artifacts, birds, mammals, and miscellaneous. The numerical statistics of each type is displayed in Table \ref{tab:dataset}. Figure \ref{fig:dataset_example} shows some examples from each big types. A full description of the dataset is shown in Figure \ref{tab:full_dataset} in the appendix. 

\begin{figure}
	 \centerline{\includegraphics[width=1\textwidth]{src/approach/dataset_examples.png}}
	 {\caption{Examples of the data classes in the dataset. The top row displays the anthropological artifacts, followed by birds in the second row, mammals in the third row, and miscellaneous in the last row.} \label{fig:dataset_example}}
\end{figure}



In order to train YOLOv5 on the detection dataset, the images are required to be accompanied by labels and masks of the objects. However, since YOLOv5 can only accept bounding boxes, the masks were converted into coordinate and size formats that conform to the YOLOv5 format (x, y, height, width). Figure \ref{fig:full_dataset_yolov5} shows a sample training batch.

\begin{figure}
	 \centerline{\includegraphics[width=1\textwidth]{src/approach/dataset_yolo.jpg}}
	 {\caption{Examples of a training batch in YOLOv5s with our dataset} \label{fig:full_dataset_yolov5}}
\end{figure}

To improve practicality and reduce training time, the image resolution for the experiment was to 256x256 pixels. This dataset was then used to train and evaluate Yolov5s for object detection and EfficientNet version B0 for classification of marine wildlife in aerial images.

\begin{table}[!h]
\centering
\begin{tabular}{lcccc}
\toprule
Type & Number of classes & Numer of samples\\
\midrule
Anthropological Artifact & 18  &  331\\
Bird & 74  & 1729\\
Mammals & 7  &  189\\
Miscellaneous & 18   & 423\\
\bottomrule
\end{tabular}
\caption{Description of the dataset according to types}
\label{tab:dataset}
\end{table}



\subsubsection{Prior knowledge of datasets}
Prior knowledge about the dataset is important for developing effective machine learning models. In the context of our marine wildlife detection and classification task, we have observed that our dataset contains center bias, meaning that the objects of interest are often located in the center of the image. This is important to keep in mind when designing our object detection models, as the model may have a tendency to focus on the center of the image and may struggle with detecting objects that are located near the edges of the image.

Another challenge is that the object of interest, i.e., the marine wildlife, is usually very small when compared to the background of the image. This makes it difficult for the model to identify the object of interest and distinguish it from the background. It is also important to note that most of the images in our dataset contain only one object, i.e., a single bird. However, there are also images where a flock of birds can be seen. This introduces another layer of complexity in object detection, as the model needs to be able to identify and differentiate between individual objects within a group. Luckily, one advantage of our dataset is that it is not affected by occlusion from clouds or bad weather or motion blurs, allowing for clearer images of the birds. 

Understanding these characteristics of the dataset can inform our selection and application of data augmentation techniques to improve the robustness of our models. For example, we may choose to apply \hyperref[CenterCrop]{CenterCrop} or \hyperref[Perspective]{Perspective} to address the center bias in our dataset, or use sharpen or CLAHE to enhance the visibility of small objects in the images. But while doing that, we have to be careful to see if the performance in pictures of a flock of bird decreases. By taking these factors into consideration, we can better prepare our models for the challenges of our specific dataset and improve their overall performance.

\subsubsection{Underpresented classes (edge cases)}
It is crucial to pay attention to the underrepresented classes (edge cases) in our dataset since they can have a significant impact on the performance of the machine learning models. In this thesis, we will compare the performance of our models before and after augmenting the dataset with GAN-generated images. We need to ensure that our evaluation includes extreme cases, where some classes have a small number of samples, such as less than 10 samples, to account for edge cases in our experiments. By doing so, we can determine if the GAN-generated images have a dramatic effect on the performance of the models, especially on underrepresented classes.  Table \ref{tab:edge-cases} shows some examples of these edge classes.

\begin{table}[ht]
\centering
\begin{tabular}{lccc}
    \toprule
    Name & Number of samples & Type \\ \midrule
        Gryllteiste  & 4 &  Bird \\
        Schnatterente & 4 &  Bird \\
        unbestimmte Larusmöwe &  4 & Bird  \\
        Schmarotzer/Spatel/Falkenraubmöwe & 5 & Bird  \\
        Wasserlinie mit Großalgen & 5 & Miscellaneous  \\
        Feldlerche & 5 & Bird  \\
        Schmarotzerraubmöwe & 5 & Bird  \\
        Grosser Brachvogel & 5 & Bird  \\
        unbestimmte Raubmöwe & 7 & Birds  \\
        Turmfalke & 7 & Bird  \\
        Trauerseeschwalbe & 7 & Bird  \\
        unbestimmter Schwan & 8 & Bird  \\
    \bottomrule
\end{tabular}
\caption{Some underpresented classes (edge classes) that will be used as input for our GAN model}
\label{tab:edge-cases} 
\end{table}

\subsubsection{Dataset stratification}
When splitting our dataset, we must use stratified sampling instead of random sampling to ensure that the proportion of each class in the training and testing sets is approximately the same as that in the entire dataset. This is because our dataset is highly imbalanced. The intuition behind this relates to the bias of most classification algorithms. They tend to weight each instance equally which means overrepresented classes get too much weight (e.g. optimizing F-measure, Accuracy or a complementary form of error) \cite{powers2015f}. By performing dataset stratification, we split the dataset 80\% into training and minimum 20\% into validation sets so they guarantees the representation of all classes in both subsets. Failure to ensure such stratification may result in the training and validation sets having different class distributions, leading to biased evaluations of the model's performance. 

